\documentclass[11pt]{article}
\usepackage[margin=1in]{geometry}
\usepackage{amsmath,amssymb,amsthm}
\usepackage[hidelinks]{hyperref}

\newtheorem{theorem}{Theorem}
\newtheorem{remark}{Remark}
\newcommand{\R}{\mathbb{R}}
\newcommand{\supp}{\operatorname{supp}}
\newcommand{\esssupp}{\operatorname{ess\mbox{-}supp}}
\newcommand{\Lip}{\operatorname{Lip}}

\title{Pure ReLU Networks Already Give Compact-Support Universality\\
\large A literature-implied full answer to the pooling question in arXiv:2204.11231}
\author{Automatic Math Discovery Run \texttt{fa\_banach\_001}}
\date{August 9, 2026}

\begin{document}
\maketitle

\section*{Source question}

Kratsios and Zamanlooy construct ReLU networks with bilinear pooling that
approximate compactly supported functions while controlling the support
\cite{KZ2022}.  On PDF page 16 they explain that replacing the exact product
by a ReLU approximation to multiplication need not preserve compact support.
They therefore leave open whether one can construct a deep ReLU network
without pooling that satisfies the hypotheses of their Lemma~2.

That lemma asks for the following.  Given a compactly supported Lipschitz map
\(f:\R^d\to\R^D\), find approximants \(f_m\) such that
\[
  \lVert f_m-f\rVert_{L^1(\R^d;\R^D)}\longrightarrow 0
\]
and such that \(\esssupp f\) and all \(\esssupp f_m\) lie in one fixed
enlarged cube.  Under these two conditions, Lemma~2 concludes that
\(f_m\to f\) in the paper's compactly-supported \(L^1\) topology \(\tau\).

\section*{Result}

Let \(\mathcal N^{\mathrm{ReLU}}_{d,D}\) denote the functions
\(\R^d\to\R^D\) realized by finite feedforward ReLU networks with affine
output layer.

\begin{theorem}[Pooling is unnecessary]
Let \(d,D\geq 1\), and let \(f:\R^d\to\R^D\) be Lipschitz with compact
essential support.  If
\[
 n_f=\min\{n\in\mathbb N:\esssupp f\subseteq[-n,n]^d\},
 \qquad Q=[-n_f-1,n_f+1]^d,
\]
then for every \(\varepsilon>0\) there is
\(N_\varepsilon\in\mathcal N^{\mathrm{ReLU}}_{d,D}\) such that
\[
 \lVert N_\varepsilon-f\rVert_{L^\infty(\R^d;\R^D)}<\varepsilon,
 \qquad \esssupp N_\varepsilon\subseteq Q.
\]
Consequently, for the cs\(L^1\)-topology \(\tau\) of arXiv:2204.11231,
\[
 \overline{\mathcal N^{\mathrm{ReLU}}_{d,D}}^{\,\tau}
 =L^1_{\mathrm{loc}}(\R^d;\R^D).
\]
\end{theorem}

\section*{Idea of the proof}

The product-and-mask route is not needed.  On a simplicial mesh of the fixed
cube \(Q\), interpolate \(f\) at the vertices.  The interpolation error is
bounded by the Lipschitz constant times the mesh diameter.  Because the cube
has a full collar outside the support of \(f\), all boundary data are zero;
the interpolant therefore extends by zero to a global compactly supported
continuous piecewise-affine map.  A pre-existing exact representation theorem
then turns this finite CPWL map into a finite ReLU network without changing it
or its support.

\section*{Proof}

For a continuous function, essential support agrees with the closure of its
nonzero set.  Indeed, if \(f(x)\neq0\), continuity makes \(f\) nonzero on a
set of positive measure near \(x\).  Hence \(f=0\) outside
\([-n_f,n_f]^d\), and in particular on a neighbourhood of \(\partial Q\).

Write \(L=\Lip(f)\).  If \(L=0\), compact support forces \(f=0\), and the
zero network proves the claim.  Assume \(L>0\).  Choose a finite conforming
simplicial triangulation \(\mathcal T_h\) of \(Q\) whose mesh diameter
\[
 h:=\max_{S\in\mathcal T_h}\operatorname{diam}(S)
\]
is smaller than \(\varepsilon/L\).  Let \(I_hf\) be the nodal affine
interpolant of \(f\) on this triangulation.  If \(x\) belongs to a simplex
\(S\) with vertices \(v_0,\dots,v_d\) and barycentric coordinates
\(\lambda_0,\dots,\lambda_d\), then
\[
 I_hf(x)=\sum_{i=0}^d\lambda_i f(v_i).
\]
Therefore
\[
 \begin{aligned}
 \lVert I_hf(x)-f(x)\rVert
 &\leq \sum_{i=0}^d\lambda_i\lVert f(v_i)-f(x)\rVert \\
 &\leq L\sum_{i=0}^d\lambda_i\lVert v_i-x\rVert
 \leq Lh<\varepsilon.
 \end{aligned}
\]

Every boundary vertex of \(Q\) has value zero.  On each boundary face the
interpolant is affine with zero vertex values, so its boundary trace is zero.
Extend \(I_hf\) by zero on \(\R^d\setminus Q\), and call the resulting map
\(g_h\).  It is globally continuous, affine on each member of a finite
polyhedral decomposition of \(\R^d\), and satisfies
\[
 \supp g_h\subseteq Q,
 \qquad \lVert g_h-f\rVert_\infty<\varepsilon.
\]
Thus every coordinate of \(g_h\) is a finite scalar CPWL function.

Theorem~5.2 on PDF page 23 of He--Li--Xu--Zheng states that every finite
scalar CPWL function on \(\R^d\) is represented exactly by a finite ReLU
DNN \cite{HLXZ2018}.  Apply that theorem to the \(D\) coordinates of
\(g_h\), run the resulting networks in parallel, and collect their outputs
with one affine output layer.  The resulting pure ReLU network
\(N_\varepsilon\) equals \(g_h\) pointwise.  It therefore has the claimed
uniform error and support.

Choose \(h_m\downarrow0\).  The networks just constructed have all their
supports in the same cube \(Q\), and
\[
 \lVert N_m-f\rVert_{L^1}
 \leq |Q|\,\lVert N_m-f\rVert_\infty\longrightarrow0.
\]
These are exactly the two hypotheses of Lemma~2 of the source paper (PDF
page 12), so \(N_m\to f\) in \(\tau\).  Lemma~1 of that paper says that
compactly supported Lipschitz maps are \(\tau\)-dense in
\(L^1_{\mathrm{loc}}\).  The usual transitivity-of-density argument used in
the source now gives the asserted \(\tau\)-density of pure ReLU networks.
\qed

\section*{Status and scope}

This completely answers the qualitative question on PDF page 16: a pure ReLU
network can meet Lemma~2, so bilinear pooling is not required for cs\(L^1\)
universality.  The result also gives uniform approximation of each compactly
supported Lipschitz target with support in the fixed one-unit enlargement
\(Q\).  It does not claim the specific width, depth, or metric-capacity bounds
of the source's quantitative Theorem~3.

The status is \texttt{literature\_implied\_answer}: the interpolation
reduction above identifies the answer, while the decisive exact CPWL-to-ReLU
representation theorem was already available before the source paper.  A
search of the run indexes and arXiv/web queries found no paper explicitly
presenting this reduction as an answer to the question in arXiv:2204.11231.

\begin{thebibliography}{9}
\bibitem{KZ2022}
A. Kratsios and B. Zamanlooy,
\emph{Do ReLU Networks Have An Edge When Approximating Compactly-Supported
Functions?}, Transactions on Machine Learning Research (2022),
arXiv:2204.11231.

\bibitem{HLXZ2018}
J. He, L. Li, J. Xu, and C. Zheng,
\emph{ReLU Deep Neural Networks and Linear Finite Elements},
arXiv:1807.03973, Theorem~5.2, PDF p.~23 (2018).
The paper records that the exact general CPWL representation theorem is due
to R. Arora, A. Basu, P. Mianjy, and A. Mukherjee,
\emph{Understanding Deep Neural Networks with Rectified Linear Units},
arXiv:1611.01491.
\end{thebibliography}

\end{document}
